\chapter{Deep learning: the autoencoder detector}
\label{ch:dl}

PCA finds the best \emph{linear} subspace for normal data. The plant is not
linear: Torricelli drains go as $\sqrt h$, the pump curve is quadratic and the
valve is equal-percentage. An autoencoder learns a \emph{non-linear} manifold
with the same scoring and the same contribution analysis. This chapter
explains the neural-network theory behind it and how it is trained here.

\section{Neural network fundamentals}

\subsection{The neuron and the layer}
A layer maps an input vector $\bm h^{(l-1)}$ to
\begin{equation}
  \bm h^{(l)} = \phi\big(W^{(l)}\bm h^{(l-1)} + \bm b^{(l)}\big),
\end{equation}
with weights $W^{(l)}$, biases $\bm b^{(l)}$ and an element-wise non-linear
activation $\phi$. Stacking layers gives a multilayer perceptron (MLP), which
by the universal approximation theorem can represent any continuous function
on a compact set given enough units.

\subsection{Activation}
pyplant uses the hyperbolic tangent,
\begin{equation}
  \tanh(a) = \frac{e^a - e^{-a}}{e^a + e^{-a}}, \qquad
  \tanh'(a) = 1 - \tanh^2(a).
\end{equation}
It is smooth, zero-centred (which suits standardised inputs) and saturates
gently. The output layer is linear, because standardised features are not
bounded.

\subsection{Learning: loss, backpropagation and Adam}
Training minimises the mean squared error over a batch $\mathcal B$,
\begin{equation}
  \mathcal L(\Theta) = \frac{1}{|\mathcal B|}\sum_{i\in\mathcal B}
  \lVert \bm z_i - f_\Theta(\bm z_i)\rVert^2 + \alpha\lVert\Theta\rVert^2 ,
\end{equation}
where the second term is a small L2 weight penalty ($\alpha = 10^{-4}$, the
scikit-learn default). Gradients $\nabla_\Theta\mathcal L$ are computed by
backpropagation, the chain rule applied layer by layer from the output back
to the input. Parameters are updated with Adam, which keeps running averages
of the gradient and its square:
\begin{align}
  \bm m_t &= \beta_1 \bm m_{t-1} + (1-\beta_1)\bm g_t, &
  \bm v_t &= \beta_2 \bm v_{t-1} + (1-\beta_2)\bm g_t^2, \\
  \Theta_t &= \Theta_{t-1} - \eta\,\frac{\hat{\bm m}_t}{\sqrt{\hat{\bm v}_t}+\epsilon},
  & \hat{\bm m}_t &= \frac{\bm m_t}{1-\beta_1^t},\;
    \hat{\bm v}_t = \frac{\bm v_t}{1-\beta_2^t},
\end{align}
with learning rate $\eta = 10^{-3}$, $\beta_1=0.9$, $\beta_2=0.999$.

\subsection{Generalisation and early stopping}
A network with enough capacity can memorise its training set, in which case it
reconstructs \emph{everything} well, faults included, and becomes useless as a
detector. 15\% of the training windows are held back as a validation set, and
training stops when the validation loss has not improved for 25 epochs
(\emph{early stopping}), up to a maximum of 1500 iterations.

\section{The autoencoder}

An autoencoder is a network trained to reproduce its own input through a
narrow \emph{bottleneck}. The encoder compresses the 48 standardised features
to 8 numbers; the decoder expands them back. Because the bottleneck is
narrower than the input, the network can reproduce only inputs that resemble
what it was trained on: it learns the manifold of normal operation.

\begin{figure}[htbp]
\centering
\scalebox{0.8}{%
\begin{tikzpicture}[x=2.3cm,y=0.42cm,
  n/.style={circle,draw=slate,fill=panelgrey,minimum size=3.2mm,inner sep=0pt},
  b/.style={n,fill=accent!40}]
  \foreach \l/\k/\style in {0/12/n,1/8/n,2/3/b,3/8/n,4/12/n}{
    \foreach \i in {1,...,\k}{
      \node[\style] (L\l-\i) at (\l,{\i-(\k+1)/2}) {};
    }
  }
  \foreach \a/\ka/\b/\kb in {0/12/1/8,1/8/2/3,2/3/3/8,3/8/4/12}{
    \foreach \i in {1,...,\ka}{\foreach \j in {1,...,\kb}{
      \draw[slate!25] (L\a-\i) -- (L\b-\j);}}
  }
  \node[below,font=\footnotesize,align=center] at (0,-7) {input $\bm z$\\48};
  \node[below,font=\footnotesize,align=center] at (1,-7) {hidden\\32, tanh};
  \node[below,font=\footnotesize,align=center] at (2,-7) {bottleneck\\8, tanh};
  \node[below,font=\footnotesize,align=center] at (3,-7) {hidden\\32, tanh};
  \node[below,font=\footnotesize,align=center] at (4,-7) {output $\hat{\bm z}$\\48, linear};
  \draw[decorate,decoration={brace,amplitude=5pt}] (-0.1,7.2) -- node[above=5pt,font=\small]{encoder} (2.1,7.2);
  \draw[decorate,decoration={brace,amplitude=5pt}] (1.9,8.6) -- node[above=5pt,font=\small]{decoder} (4.1,8.6);
\end{tikzpicture}}
\caption{The 48--32--8--32--48 autoencoder (drawn with fewer units). It is
trained to output its input.}
\label{fig:ae}
\end{figure}

\begin{center}
\begin{tabular}{@{}lll@{}}
\toprule
Setting & Value & Field in \code{AutoencoderDetector} \\
\midrule
hidden layers & (32, 8, 32) & \code{hidden} \\
activation & tanh & fixed \\
optimiser & Adam, $\eta = 10^{-3}$ & fixed \\
maximum iterations & 1500 & \code{max_iter} \\
early stopping & 15\% validation, patience 25 & fixed \\
random seed & 0 & \code{seed} \\
\bottomrule
\end{tabular}
\end{center}

\subsection{Relationship to PCA}
A linear autoencoder (identity activations) trained with MSE learns exactly
the principal subspace: its optimum spans the same space as the top $k$
principal components. PCA is therefore the linear special case, and the
autoencoder earns its place only if the non-linearity buys detection that PCA
misses. That is why both are trained and the comparison is kept on record.

\section{Scoring, threshold and isolation}

The autoencoder uses the same machinery as PCA:
\begin{equation}
  a_\text{AE}(\bm x) = \lVert \bm z - f_\Theta(\bm z) \rVert^2,
\end{equation}
the same calibration quantile, the same health index and the same
per-channel contributions of the residual (Chapter~\ref{ch:ml}). Only the
reconstruction changes.

\begin{figure}[htbp]
\centering
\scalebox{0.8}{%
\begin{tikzpicture}[node distance=7mm and 10mm]
  \node[box] (x) {window $\bm x$};
  \node[box,right=of x] (z) {standardise\\$\bm z$};
  \node[hot,right=of z] (f) {reconstruct\\$\hat{\bm z}=f(\bm z)$};
  \node[box,right=of f] (r) {residual\\$\bm r=\bm z-\hat{\bm z}$};
  \node[box,above right=4mm and 10mm of r] (s) {score $\lVert\bm r\rVert^2$\\$\to$ HI, flag};
  \node[box,below right=4mm and 10mm of r] (c) {$r_j^2$ by channel\\$\to$ top channel};
  \draw[flow] (x) -- (z);
  \draw[flow] (z) -- (f);
  \draw[flow] (f) -- (r);
  \draw[flow] (r) |- (s);
  \draw[flow] (r) |- (c);
\end{tikzpicture}}
\caption{One scoring pass, identical for PCA and the autoencoder.}
\label{fig:aescore}
\end{figure}

\section{Practical advice}

\begin{itemize}
  \item \textbf{More data before more layers.} With the Quick preset there are
        only a few normal runs; a bigger network will memorise them. Use the
        Standard or Thorough preset before changing \code{hidden}.
  \item \textbf{Bottleneck width} controls the trade-off: too wide and faults
        reconstruct well (missed detections), too narrow and normal variation
        does not (false alarms). 8 is a good starting point for 48 features.
  \item \textbf{Seeds matter.} Training is stochastic. Compare methods over
        two or three seeds before concluding one is better.
  \item \textbf{Check against PCA.} If the autoencoder does not beat PCA on
        detection rate and false alarms, deploy PCA. The selection rule does
        this automatically.
\end{itemize}

To change the architecture:
\begin{lstlisting}
from pyplant.ml.detector import AutoencoderDetector
det = AutoencoderDetector(quantile=0.995, hidden=(64, 16, 64), max_iter=3000, seed=1)
\end{lstlisting}
